\section{Introduction}
\label{sec:introduction}

Sixth-generation (6G) wireless networks are expected to provide ubiquitous
connectivity for a massive number of Internet-of-Things (IoT) devices
\cite{cai2026,xu2026}. To serve devices in remote areas with limited terrestrial
coverage, low Earth orbit (LEO) satellites complement cellular networks for
applications such as smart agriculture and climate monitoring \cite{shen2025,kodheli2021}. These
power-limited devices transmit short packets sporadically, with only a small
fraction active at a given time \cite{shen2025,cai2026,mei2022}. In grant-based random
access, handshaking before transmission incurs significant
signaling overhead, while the long propagation delay of terrestrial-satellite links
further increases access latency \cite{cai2026,xu2026}. Accordingly, grant-free
random access (GFRA) allows active devices to transmit pilots and data directly
without first obtaining a grant, reducing signaling overhead and access latency in LEO satellite IoT
systems \cite{cai2026,xu2026}.

Supporting GFRA over LEO links requires handling large delays and Doppler shifts
caused by long propagation distances and high satellite mobility. The 3rd Generation
Partnership Project (3GPP) Release 17 framework addresses these effects through
device-side precompensation assisted by the global navigation satellite system
(GNSS) \cite{3gpp38300}. With such precompensation, many existing satellite GFRA
schemes assume perfect compensation or small residual delays and Doppler shifts at
the receiver \cite{shen2025}. Frequent estimation and compensation increase
energy consumption and reduce battery lifetime, making these schemes less suitable
for power-limited remote IoT devices \cite{shen2025}. This motivates reducing
dependence on GNSS, consistent with the GNSS-independent operation considered in the
Release 19 study on satellite access \cite{3gpp22865}. With reduced reliance on
device-side precompensation, the satellite receiver must accommodate
large differential delays and Doppler shifts, which can cause severe inter-symbol
and inter-carrier interference in conventional orthogonal frequency division
multiplexing (OFDM) systems \cite{shen2025,huang2026}.

Orthogonal time frequency space (OTFS) modulation exploits the quasi-static
properties of channels in the delay--Doppler domain to mitigate these effects
\cite{shen2025,hadani2017}. In \cite{shen2025}, multi-frame OTFS transmission
handles large differential delays and Doppler shifts at the satellite receiver,
allowing devices to operate without GNSS assistance or precompensation of these
offsets. Affine frequency division multiplexing (AFDM)
offers a low-complexity modulation structure based on one-dimensional chirp
signals \cite{ding2026,huang2026}. Its modulator can be implemented using a
single inverse fast Fourier transform and two stages of pointwise chirp
multiplication \cite[Secs.~II-B and III-A]{bemani2022}.

Accurate joint activity detection and channel estimation (JADCE) is
essential for AFDM-GFRA \cite{ding2026}. Existing
methods exploit sporadic device activity to formulate JADCE as a compressed sensing
(CS) problem and incorporate structured priors to capture channel sparsity
\cite{shen2025,ding2026,ke2020}. The scheme in \cite{ding2026} further uses a graph neural
network to adaptively learn sparsity patterns, while retaining a predefined
delay--Doppler grid for the sparse channel representation. In practical
propagation environments, significant power leakage caused by fractional Doppler
shifts and fractional time delays can weaken the sparsity of the effective channel
matrix \cite{huang2026}. This motivates learning effective-channel
priors from data to capture channel structures beyond the assumed sparse
representation \cite{cai2026,huang2026}.

Score-based generative models learn channel priors from samples by estimating the
score function of the channel distribution \cite{song2019,arvinte2023}. In
\cite{huang2026}, score-based learning is performed directly on the AFDM effective
channel to address fractional delay--Doppler off-grid leakage in channel estimation.
For massive random access, \cite{ying2026} proposes a two-stage algorithm that first
identifies active devices through a Transformer-based detector and then uses
generative diffusion models for joint channel estimation and data detection.
Both approaches rely on iterative posterior sampling, which requires
repeated score function evaluations and incurs substantial computational cost
\cite{huang2026,ying2026}.

The score-based turbo message passing (STMP) algorithm in \cite{cai2026}
incorporates generative channel priors for JADCE without iterative posterior
sampling. Through Tweedie's formulas, separate first- and second-order score networks provide
posterior means and variances for single-step minimum mean-squared error (MMSE)
denoising under equivalent additive white Gaussian noise (AWGN) observations.
The relation between denoiser derivatives and estimation uncertainty
\cite{dytso2023} offers another way to obtain variance information:
a trained denoiser can also quantify uncertainty without
additional training \cite{manor2024}.
This allows the learned channel prior to supply both channel estimates
and variance information for message passing through the same denoiser.
Moreover, STMP computes activity probabilities using Gaussian extrinsic messages
obtained by projecting the channel posterior onto a Gaussian distribution
\cite{cai2026}. For a non-Gaussian channel prior, matching the posterior mean and
variance does not generally preserve the marginal likelihood. The relation between
density ratios and class probabilities \cite{menon2016} offers an alternative:
estimating the activity likelihood ratio directly from samples, without relying
on this Gaussian approximation or explicitly evaluating the marginal likelihood.

Against this background, we investigate GNSS-independent AFDM-GFRA for
LEO satellite IoT systems. Our main contributions are summarized as follows:

\begin{enumerate}[label=\arabic*)]
\item We propose a multi-frame AFDM transmission scheme that accommodates large
differential delays and Doppler shifts at the satellite receiver without
requiring device-side precompensation. We derive the input-output relationship
and show that the effective channels across frames differ only by
Doppler-induced phase rotations. Separating these phase rotations allows the
received pilot frames to be jointly described by a shared effective channel
matrix. Additional pilot frames therefore provide more observations without
increasing the number of channel coefficients to be estimated. Compared with
the multi-frame OTFS representation of Shen et al.~\cite{shen2025}, this
effective-channel representation contains fewer coefficients and provides a
compact model for learning the channel prior.

\item We learn a score-based generative prior for AFDM effective channels without
imposing sparse or grid-aligned channel assumptions. Building on Tweedie's
formula and the relation between denoiser derivatives and estimation
uncertainty, we develop a channel denoiser that provides both channel estimates
and average variance without requiring a separately trained variance model.

\item We develop a message-passing JADCE receiver comprising two modules:
multi-frame linear estimation (MFLE), which jointly processes the pilot frames
to suppress inter-user interference, and generative-prior-aided joint denoising
(GPJD), which incorporates the proposed channel denoiser. Exploiting the relation
between density ratios and class probabilities, we incorporate a learned
density-ratio estimator into GPJD, enabling activity detection under the learned
channel prior without explicit marginal-likelihood evaluation. The two modules
iteratively exchange extrinsic messages to perform JADCE. Simulations show that
the proposed receiver outperforms the STMP baseline adapted to AFDM in activity
detection and channel estimation, and requires fewer pilot frames to achieve
the same channel-estimation accuracy.
\end{enumerate}

The remainder of this paper is organized as follows. Section~\ref{sec:system}
presents the system model and the multi-frame AFDM input--output relationship.
Sections~\ref{sec:algorithm}--\ref{sec:overall} describe the turbo message-passing
framework, the learned channel denoiser and activity classifier, and the overall algorithm with its computational
complexity analysis, respectively. Section~\ref{sec:simulations} presents the
simulation results. Section~\ref{sec:conclusion} concludes this paper.
